% paper/sec00-example-overview.tex -- FORMAT EXAMPLE for the reading of the
% paper's introduction, not a section of any book.
%
% The first chapter of Part II reads the introduction and, in doing so, gives
% the map of the paper: what the main theorem says and how its proof runs
% (CLAUDE.md, Phase 5, "the reading of the introduction").  It comes FIRST in
% the reading order and is WRITTEN LAST, after every other Part II chapter has
% been approved, so that each step of the map can point at the chapter that
% proves it.  Its chapter number is reserved in main.tex from the start.
%
% Copy this file to paper/secNN-SLUG.tex, NN being the number of the
% introduction in the paper (00 if it is section 0 or unnumbered).
%
% Notation and definitions are the one place where this chapter departs from
% the writing guide (PHASE45_WRITING_GUIDE.md, section 2.9):
%   - everything it uses from Part I is recalled at the start, in one place;
%   - anything defined in another Part II chapter is not defined here: at its
%     first use it gets one phrase for its meaning, the paper's location by
%     \cite, and the book's location by \ref.
% The exception ends with this chapter.  No other chapter leans on it.
%
% main.tex deliberately does NOT \include it.
%
% ASCII only under tex/.
\chapter{<Title: the main theorem and the shape of its proof>}
\chaptermark{<Short running head>}
\label{ch:read-intro}

\begin{chapterheader}
\textbf{Purpose.}
To read \S<n> of the paper, and with it the map of the whole paper: what the
main theorem says, and the steps of its proof, each with the chapter of this
part where it is proved.

\medskip
\textbf{Assumed.}
Chapter~\ref{ch:SLUG} for <what>; <further preparatory chapters, each with
what it supplies>.  Objects that the paper defines in its later sections are
used here before their definition, each with a pointer to the chapter of this
part that defines it.

\medskip
\textbf{Supports.}
\S<n> of the paper, pp.~<a>--<b>.  The chapters that follow in this part each
carry out one or more steps of the map in
Section~\ref{sec:read-intro-map}.

\medskip
\textbf{How far things are proved.}
The map proves nothing; each step names the chapter that proves it.  <Say
whose the steps are: the paper's own outline, unfolded, or ours -- if the
introduction gives no outline, or the steps are cut differently from it, say
so here.>
\end{chapterheader}

% The recall of Part I: every symbol and notion of Part I used anywhere in this
% chapter, in this one paragraph, however many sections use it.  The form of a
% recall is unchanged -- one phrase for the meaning and a \ref.  Symbols fixed
% throughout the book are not repeated.
\noindent\textbf{From the preparatory chapters.}
<$X$ denotes <one phrase>, as in Definition~\ref{def:SLUG}; <...>.>

%% =====================================================================
\section{<The main theorem>}
\label{sec:read-intro-statement}

% Sections follow the order of the introduction.  Its other parts --
% motivation, earlier results, the paper's comments on its method -- are read
% in sections of their own, as any section of the paper is read.

<A symbol or a notion defined in another chapter of this part is cited at its
first use -- just before the statement, if the statement needs it: a
\emph{<notion>} is <one phrase for its meaning>
(\cite[Def.~<n.m>]{<PAPER>}; Definition~\ref{def:read-SLUG}).>

\begin{thm}[{\cite[Thm.~<n.m>]{<PAPER>}}]
\label{thm:read-intro-main}
<The main theorem, stated exactly, in the notation of the chapter that proves
it.  It is proved in Chapter~\ref{ch:read-SLUG}.>
\end{thm}

<If the paper has several main results, state each, and say which follows from
which.>

%% =====================================================================
\section{<The shape of the proof>}
\label{sec:read-intro-map}

% The map is a numbered list, outside any proof environment, and carries no
% verification mark: nothing is proved here (PHASE45_WRITING_GUIDE.md,
% sections 5 and 6).  For each step, the geometric meaning comes first --
% what is built, what it guarantees, and what that lets the next step do --
% written for a reader who opens the paper for the first time.

<One paragraph on the strategy as a whole, before the steps.>

\begin{enumerate}
\item \emph{<What the step builds, or what it reduces to.>}
  <Its geometric meaning: what is built and what it guarantees.  Why it
  holds, in one paragraph.>
  Proved in \cite[\S<k>]{<PAPER>}; Chapter~\ref{ch:read-SLUG}.
\item \emph{<The next step.>}
  <...>
\end{enumerate}

%% =====================================================================
\section{Summary: where each statement is proved}
\label{sec:read-intro-summary}

% The third column points at where each item is PROVED: the items of this
% chapter are the statements that the rest of Part II proves.
\begin{center}
\small
\begin{tabular}{@{}L{4.6cm}lL{5.9cm}@{}}
\toprule
Item & Here & Proved in \\
\midrule
<the main theorem>
  & \ref{thm:read-intro-main}
  & Chapter~\ref{ch:read-SLUG}; \cite[\S<k>]{<PAPER>} \\
\bottomrule
\end{tabular}
\end{center}
